% =====================================================================
\section{Сон и энергосбережение}
\label{les:sleep}
% =====================================================================
\lessonintro
  {разобраться в режимах сна ESP32-S3, научиться будить чип таймером и кнопкой, посчитать, сколько проработает батарейка}
  {урок~\ref{les:gpio-in} (кнопка), урок~\ref{les:adc} (измерение напряжения), урок~\ref{les:radio} (Wi-Fi как главный потребитель)}
  {«датчик на батарейке»: просыпается раз в 30 секунд или от кнопки, измеряет и снова засыпает}
  {кнопка и светодиод (рис.~\ref{fig:button}), потенциометр (рис.~\ref{fig:pot}), мультиметр (по желанию)}

\subsection{Зачем устройству спать}

Метеостанция на даче, датчик протечки под ванной, умный ошейник для кота --- все они работают от батареек
месяцами. Секрет прост: почти всё время они \textbf{спят}, а просыпаются на секунду --- измерить, отправить
данные --- и снова засыпают. Работающая ESP32-S3 с Wi-Fi потребляет в десятки тысяч раз больше, чем спящая.

\begin{figure}[H]
\centering
\begin{tikzpicture}
\begin{axis}[
  xbar, width=0.8\textwidth, height=5.8cm,
  xmode=log, log basis x=10, log origin x=infty,
  xmin=0.002, xmax=20000,
  xlabel={потребление, мА (каждое деление --- в 10 раз больше)},
  symbolic y coords={Hibernation,Deep-sleep,Light-sleep,Работа без Wi-Fi,Работа + Wi-Fi},
  ytick=data, y tick label style={font=\sffamily\small},
  tick label style={font=\scriptsize},
  nodes near coords, nodes near coords align={horizontal},
  every node near coord/.append style={font=\scriptsize,anchor=west,xshift=1pt},
  point meta=explicit symbolic,
  bar width=5mm, enlarge y limits=0.15,
  xmajorgrids, grid style={gray!25},
]
\addplot[fill=esp!70,draw=espdark] coordinates {
  (0.007,Hibernation) [$\approx$ 7 мкА]
  (0.008,Deep-sleep) [$\approx$ 8 мкА]
  (0.24,Light-sleep) [$\approx$ 240 мкА]
  (30,Работа без Wi-Fi) [$\approx$ 30--50 мА]
  (300,Работа + Wi-Fi) [$\approx$ 100--300 мА]
};
\end{axis}
\end{tikzpicture}
\caption{Сколько потребляет сам чип в разных режимах (примерно, по документации).
Отладочная плата с преобразователем питания и светодиодом питания потребляет заметно больше}
\end{figure}

\begin{center}\small
\begin{tabularx}{\textwidth}{@{}l X X@{}}
\toprule
Режим & Что работает & Что при пробуждении \\
\midrule
Light-sleep (лёгкий сон) & всё «заморожено», память сохранена & программа продолжается со следующей строки \\
Deep-sleep (глубокий сон) & только часы и маленькая память RTC & \textbf{перезапуск}: \code{setup()} начинается заново \\
Hibernation (спячка) & только часы & перезапуск \\
\bottomrule
\end{tabularx}
\end{center}

Главное о глубоком сне: при пробуждении программа начинается сначала, а обычные переменные забываются.
Сохранить данные можно в маленькой памяти RTC --- для этого переменную помечают \code{RTC\_DATA\_ATTR}.

\subsection{Сколько проработает батарейка}

Типичный «датчик на батарейке» живёт по такому расписанию:

\begin{figure}[H]
\centering
\begin{tikzpicture}[node distance=6mm]
  \node[blk,fill=blkrtc] (wake) {Проснуться};
  \node[blk,fill=blkper,right=of wake] (meas) {Измерить};
  \node[blk,fill=blkrf,right=of meas] (wifi) {Отправить\\ по Wi-Fi};
  \node[blk,fill=blkcpu,right=of wifi] (sleep) {Спать\\ 10 минут};
  \draw[arr] (wake) -- (meas); \draw[arr] (meas) -- (wifi); \draw[arr] (wifi) -- (sleep);
  \draw[arr] (sleep.south) -- ++(0,-6mm) -| node[pos=0.25,below,font=\scriptsize]{будильник-таймер} (wake.south);
\end{tikzpicture}

\vspace{4mm}
\begin{tikzpicture}[x=0.9cm,y=1cm]
  \draw[->] (0,0) -- (15.5,0) node[right,font=\scriptsize]{$t$};
  \draw[->] (0,0) -- (0,2.4) node[above,font=\scriptsize]{ток};
  \draw[very thick,esp] (0,0.05) -- (1,0.05) -- (1,1.9) -- (1.7,1.9) -- (1.7,0.05) -- (7.5,0.05) -- (7.5,1.9) -- (8.2,1.9) -- (8.2,0.05) -- (14,0.05) -- (14,1.9) -- (14.7,1.9) -- (14.7,0.05) -- (15.2,0.05);
  \node[font=\scriptsize,align=center] at (4.4,0.6) {сон: микроамперы};
  \node[font=\scriptsize,align=center,anchor=south] at (1.35,1.95) {работа: 1--3 с,\\ 100+ мА};
\end{tikzpicture}
\caption{Устройство почти всё время спит. Похоже на ШИМ из урока~\ref{les:pwm} с очень маленьким коэффициентом заполнения!}
\end{figure}

Средний ток считается так же, как среднее напряжение ШИМ:
\[
  I_\text{ср} = \frac{I_\text{работа}\cdot t_\text{работа} + I_\text{сон}\cdot t_\text{сон}}{t_\text{работа}+t_\text{сон}}
  = \frac{120~\text{мА}\cdot 2~\text{с} + 0{,}01~\text{мА}\cdot 598~\text{с}}{600~\text{с}}
  \approx 0{,}41~\text{мА}.
\]
Аккумулятор ёмкостью 2000~мА·ч прослужит $2000 / 0{,}41 \approx 4900$ часов --- около \textbf{полугода}.
А без сна, при 120~мА, --- меньше суток!

\subsection{Сон по будильнику-таймеру}

\codecap{Пробуждение по таймеру со счётчиком в RTC-памяти}
\codeinput{l14_timer_sleep}

\begin{caution}
При подключении через разъём USB плата во сне «исчезает» из компьютера, и Serial Monitor отключается.
Это нормально. Если новая прошивка не загружается, войди в режим прошивки вручную (BOOT + RST, урок~\ref{les:ide}).
\end{caution}

\subsection{Пробуждение кнопкой}

Разбудить чип может и уровень на выводе (подходят \pin{0}--\pin{21}). Кнопка на \pin{5} подходит. Во сне
обычная подтяжка из урока~\ref{les:gpio-in} отключается, поэтому включаем специальную подтяжку, которая
работает и во сне:

\codeinput{l14_ext0}

\subsection{Практика: датчик на батарейке}

Соберём всё вместе. Используем кнопку и светодиод (рис.~\ref{fig:button}) и потенциометр вместо
«настоящего» датчика (рис.~\ref{fig:pot}). Плата просыпается раз в 30 секунд \emph{или} от кнопки,
измеряет напряжение, коротко мигает светодиодом и снова засыпает. Счётчики пробуждений хранятся в RTC-памяти.

\begin{figure}[H]
\centering
\begin{tikzpicture}[node distance=7mm and 9mm]
  \node[blk,fill=blkrtc] (wk) {Пробуждение};
  \node[blk,fill=blkper,right=of wk] (why) {Кто разбудил?\\ \scriptsize\code{esp\_sleep\_get\_wakeup\_cause()}};
  \node[blk,fill=blkrf,right=of why] (work) {Измерить,\\ мигнуть, напечатать};
  \node[blk,fill=blkcpu,below=of work] (sl) {Deep-sleep};
  \node[blk,fill=blkmem,below=of wk,align=center] (t) {таймер 30 с\\ или кнопка};
  \draw[arr] (wk) -- (why); \draw[arr] (why) -- (work); \draw[arr] (work) -- (sl);
  \draw[arr] (sl) -- (t); \draw[arr] (t) -- (wk);
\end{tikzpicture}
\caption{Жизненный цикл «датчика на батарейке»}
\end{figure}

\codecap{Датчик на батарейке}
\codeinput{l14_sensor_node}

\begin{summary}
\begin{itemize}[leftmargin=*]
  \item Глубокий сон снижает ток до микроампер, но пробуждение --- это перезапуск программы.
  \item Нужные данные храним в \code{RTC\_DATA\_ATTR}, причину пробуждения узнаём через \code{esp\_sleep\_get\_wakeup\_cause()}.
  \item Будить можно таймером, кнопкой, сенсорным входом.
  \item Время жизни от батареи определяется долей времени, которую устройство не спит.
\end{itemize}
\end{summary}

\begin{quiz}
\begin{enumerate}
  \item Что произойдёт с обычной переменной после глубокого сна?
  \item Во сколько раз работающая с Wi-Fi плата потребляет больше спящей?
  \item Как изменится время работы батарейки, если просыпаться раз в минуту, а не раз в 10 минут?
\end{enumerate}
\end{quiz}

\begin{tasks}
\begin{enumerate}
  \item Измерь ток платы мультиметром в работе и во сне.
  \item Добавь отправку измерения на веб-сервер или в Serial по Wi-Fi и посчитай, сколько проработает батарейка.
  \item Показывай число пробуждений на OLED-экране (урок~\ref{les:i2c}) и гаси экран перед сном.
\end{enumerate}
\end{tasks}
